\subsection{Дописывание получения движений токена по workflow}

\subsubsection{Задача}
\begin{itemize}
  \item Получение истории и статуса токенов переведено на асинхронную схему: Camunda самостоятельно публикует уведомления о событиях в брокер сообщений. 
  \item На стороне бизнес-логики реализован модуль подписки на эти уведомления без периодического опроса (polling).
  \item Подключение брокера выполнено в виде плагина Camunda Engine, активируемого при старте standalone-сборки. 
  \item Система корректно восстанавливает состояние подписки после перезапуска.
\end{itemize}

\subsubsection{Аналитика задачи}

Задача была затрона на обсуждении командных задачек 05.03.2026. 
Изначально эта часть системы была сделана с минимальными временными затратами,
поэтому легкость решения влияла на загрузку системы прямолинейно. 
В начале в имеющийся механизм меня начал погружать руководитель группы, так как нужно было дать вводный контекст по коду.
Основная проблема заключалась в том, что отсутвовал событийный механизм оповещения передвижения токена по бизнес процессу. 
Поэтому самым быстрый вариантом было опрашивать каждый заданный период времени камунду, искать по истории место токена. 
Для быстрого отклика системы период времени был сделан очень маленьким -- 0.5 секунды. Действия, которые происходили каждые 0.5 секунд, 
были достаточно трудоёмкие и излишне перегружали системы. 
Даже при таком маленьком инервале опроса, реактивность оповещения пользователя иногда отставала.
Во время обсуждения текущей ситуации с нами присутствовал системный аналитик, который параллельно формулировал задачу для разрабоки.
Далее задача была передана для согласования архитектору, который перевёл её далее на меня, а я перевёл её в статус "В процессе".
Архитектор добавил в задачу пожелание: убрать запись исторических записей в основную базу данных. Сама история оказалась бесполезной для камунды, 
она нужна только для отслеживания процесса в бизнес логике или аудита событий.

\subsubsection{Исследование точек внедрения}
После того, как я взял задачу в выполнение, сразу решил изучить внутренние механизмы основного модуля камунды,
который и служит центровой точкой для любых манипуляций с процессом. По коду и документации [3] стало видно,
что имеет смысл взять $HistoryEventHandler$, который является абстрактным классом и предоставляет сигнатуры нужных нам методов.
Архитектура будет выстроена следующим образом -- над основным сервисом $Camunda$ будет подключаться дополнительное расширение,
которое будет содержать сам $Handler$, который будет отвечать за пересылку сообщений от общей точки получения сообщений. 
Сама система уведомлений в двигателе $camunda$ устроена следующим образом: 
всем объектам, которые реализуют интерфейс HistoryEventHandler и переопределяют его методы,
будет подаваться сигнал о сообщении и само сообщение в метод $handleEvent$/$handleEvents$.
Далее в плагине я решил предусмотреть так же некую обработку сообщений перед пересылкой её в kafka. 
Обработка заключается в упаковке полученной сущности в объект для транспортировки. Далее сообщение будет передаваться 
в топик $kafka$ и после этого на стороне бизнес логики будет обработчик сообщений от $kafka$, 
который будет брать сообщение и далее обрабатывать. На уровне бизнес логики обработка сообщений 
будет сводиться к обработке трёх видов событий. Так и будет происходить синхронизация процесса из 
$Camunda$ с состоянием того же процесса в бизнес логике. События будут следующего вида: процесс завершился/отменился, тогда мы должны изменить его статус в бизнес логике; 
задача в процессе создана/удалена, тогда создаём/закрываем шаг; активность сменилась, тогда двигаем lifecycle-состояние. При этом нужно предусмотреть механизм повторной обработки при ошибке.

\subsubsection{Handler}

Я использовал достаточно простую идею для написания этого сервиса на сторона бизнес логики: приходит событие из $kafka$, я его забираю и отражаю в нашей модели. Далее как раз логика разделяется на 3 вида событий.

Начал я с метода $consumeProcessInstanceEvent$ — самый простой из трёх. Ищу объект по             
$rootProcessInstanceId$. Когда нашёл — смотрю в каком состоянии завершился      
процесс: если «прерван» или «прерван извне» — ставлю «отменён». Если «завершён»
— всё прошло по сценарию и тогда можно поставить «финиш». 
Заодно, пока уже есть объект в быстром доступе, фиксирую дату завершения, если она ещё не проставлена.
После этого отдаю в сервис $workflowObjectSubscriberBean$ в метод $workflowFinished$,
которые ранее были уже реализованы, так как наша задача не переписать с нуля, а внедрить новый механизм.                                                                         
Следующий случай — когда объект вовсе не нашёлся. Это меня попросили считать не ошибкой: событие просто
прилетело раньше, чем объект успел создаться, такое может быть из-за задержек сети и других накладных расходов.
Таких гонок не избежать, поэтому я просто возвращаю то же событие обратно в очередь с задержкой в 100
миллисекунд. Оно обождёт и попробует снова. 

Следующий вариант события я выделил в $consumeTaskInstanceEvent$ — этот случай немного сложнее.
Здесь мне достаточно базового объекта ($BaseWorkflowObject$), без лишней загрузки сторонних данных из медленной персистенции.
На «удалить» и «завершить» логика одинакова — задача закрыта в контексте процесса, замечу это не значит что закрыт процесс, передаю в $processFinishedTask$, 
который был реализован ранее и у нас опять получилось встроиться в текущий механизм без его переписывания.
На «создать» — другой случай. Тут уже нужен полный объект с данными из персистенции, поэтому сначала вызываю $createWorkflowObjectFromNuxeo$, и только потом иду в   
$processCreateTaskInstanceEvent$. Там по $taskDefinitionKey$ нахожу шаблон шага и передаю всё в $processIncomingTask$.                                             
Был старый путь, который откладывал обработку в зависимости от текущего статуса объекта. Я его закомментировал: сейчас логика упрощена, всегда идём напрямую, но нужно сохранить логику старой реализации до окончания тестирования.

И последний сценарий ($consumeActivityInstanceEvent$) — отвечает за переходы жизненного цикла. 
Здесь я внедрил фильтр: меня интересует только тип $intermediateMessageCatch$. Остальные 
типы — шлюзы, подпроцессы, служебные задачи — тоже генерируют события, но они мне не нужны.
Дальше собираю строку состояния: беру префикс по типу активности и прибавляю $activityId$. 
Если в идентификаторе есть символ \# — обрезаю всё после него, это служебный суффикс Камунды, в моей модели он не нужен.                                                                               
Важно проверить перед тем как что-то менять — не стоит ли объект уже в этом состоянии и если стоит — выхожу без лишних действий, если переход разрешён
жизненным циклом — достаю $LifecycleTransition$, из него беру состояние и передаю в $workflowMoveToTransient$, 
где старая логика сделает перевод объекта в некоторое техническое состояние,
при котором действия пользователя с объектом будут невозможны. Объект движется дальше по схеме.

Код обработчика можно найти в \textbf{Приложении Б}.

\subsubsection{Сервисы для обработки событий на уровне плагина}

Плагин меня попросили создать в отдельном проекте и через сборку контейнера внести в плагин в другой контейнер. Данный плагин должен
перехватывать события от камунды в режиме реального времени и публиковал их в шину
данных Kafka, откуда любой заинтересованный сервис мог бы их получить, в частности наш перехватчик событий на уровне бизнес логики.

Прежде чем приступить к реализации, я изучил документацию Camunda BPM по       
механизмам расширения движка [4]. Camunda предоставляет интерфейс
ProcessEnginePlugin, позволяющий вмешаться в процесс инициализации движка, а   
также возможность добавить стандартный обработчик истории через метод
setHistoryEventHandler. Как я ранее проанализировал -- интерфейс HistoryEventHandler, перехватывает абсолютно
все события движка централизованно, независимо от конкретного процесса. Именно
этот подход и был выбран как основной. 

Проект я разбил на три Maven-модуля, чтобы разделить ответственности.
                                                                                 
Модуль $event-listener-lib$ — общая библиотека, содержащая DTO-классы событий,   
перечисления, конфигурацию Kafka-продюсера. Этот модуль я решил подключать как зависимость и в плагин, и в клиентскую часть.             
Модуль $event-listener$ — собственно плагин для Camunda, регистрируемый в движке 
и перехватывающий события истории. Модуль $event-listener-client$ — клиентская библиотека для любого сервиса,       
желающего подписаться на события Camunda, этот модуль было решено сделать заранее, так как у архитектора были ещё планы на интеграции данных событий.                        
  
Первым делом я написал DTO-классы для передачи данных о событиях. 
Базовый абстрактный класс BaseEvent содержит поля, общие для всех
типов событий: тип события, класс события, идентификаторы корневого и текущего 
экземпляра процесса, ключ определения процесса и время завершения. 
От BaseEvent я отнаследовал три конкретных класса. $ActivityInstanceEvent$          
дополнительно содержит идентификатор и тип активности. 
$TaskInstanceEvent$ содержит идентификатор задачи и ключ её        
определения. $ProcessInstanceEvent$ добавляет
поле состояния процесса. Для классификации событий были введены два перечисления: $EventClass$ -- \texttt{PROCESS\_INSTANCE}, \texttt{ACTIVITY\_INSTANCE}
или \texttt{TASK\_INSTANCE} и $EventType$, который определяет, что   
именно произошло: \texttt{START}, \texttt{END}, \texttt{CREATE}, \texttt{DELETE}, \texttt{COMPLETE}. Аналогичным образом я написал $ProcessInstanceState$ —
перечисление всех возможных состояний процесса.
Параметры подключения к брокеру и название топика вынесены в переменные        
окружения через вспомогательный класс $SystemEnvironmentUtil$. 
Это позволяет деплоить плагин в разных окружениях без пересборки.
Так же для конфига я предусмотрел стандартные значения, чтобы плагин мог запускаться без явного указания конфигурации.
Конфигурация $Kafka$-продюсера регистрирует фабрику, шаблон и
сериализатор. Я ввёл собственное поле $eventClass$, по которому консюмер самостоятельно определяет
нужный конкретный тип при десериализации. Сам продюсер EventProducer принимает $BaseEvent$ и  
отправляет его в настроенный топик.

После этого я начал писать главный модуль -- event-listener.                                
Точка входа плагина — класс $EventListenerPlugin$, реализующий интерфейс         
$ProcessEnginePlugin$. В методе вызываемом Camunda до полного запуска движка, поднимается собственный контекст плагина, который и помог мне быстро переиспользовать
все инструменты $Spring-boot$. После запуска контекста из него
извлекается бин $HistoryEventHandler$ и устанавливается в конфигурацию движка.
Класс EventHandlerKafka реализует интерфейс HistoryEventHandler, как раз он и является
получателем всех событий истории от движка.
Его логику я специально сделал минималистично -- передать на преобразование и, если результат есть, отправить в Kafka. 
Метод $handleEvents$ делегирует $handleEvent$, чтобы не
дублировать логику. Вся работа по фильтрации и маппингу сосредоточена в классе $EventHandlerManager$.
Изначально я забыл вписать фильтрацию процессов и демонстрационные процессы засоряли топик служебными событиями, но далее был дописан фильтр.                                                                   
Далее я проверяю тип события и создаю объект для пересылки. Логика фильтрации по типам событий строилась постепенно — в ходе тестирования я несколько раз запускал тестовые
процессы и анализировал поток событий, определяя, какие из них несут реальную
смысловую нагрузку: для $ActivityInstance$ публикуется только событие $start$, и только если         
активность не является userTask — задачи передаются отдельным каналом через $TaskInstance$;
для $TaskInstance$ публикуются события $create$, $complete$ и $delete$;
для $ProcessInstance$ отсекаются служебные события $update$, $create$, $start$,      
$migrate$, а публикуются лишь события завершения и отмены процесса.

Клиентская библиотека лежит в отдельном модуле $event-listener-client$.
В ходе обсуждения с архитектором было решено реализовать не
просто $Kafka$-консюмер, а готовую библиотеку, которая реализует интерфейс и получает уже
типизированные объекты событий, не занимаясь самостоятельно десериализацией.

Интерфейс $EventListenerHandler$ определяет три метода. Каждый метод
возвращает объект будущего, что даст потребителю возможность
обрабатывать события асинхронно. $BaseEvent$ является абстрактным классом, стандартный
$Jackson$-десериализатор не может самостоятельно определить, в какой
конкретный тип преобразовать JSON, поэтому я реализовал
$BaseEventDeserializer$.

Конфигурация консюмера KafkaConsumerConfig намеренно выставляет
concurrency = 1 и ручное подтверждение offset через режим
\texttt{MANUAL\_IMMEDIATE}. Это я сделал чтобы обеспечить упорядоченную обработку событий. 
В ходе промежуточных тестов я столкнулся с проблемой --
один бизнес-процесс может порождать несколько событий практически одновременно, и если
обрабатывать их в произвольном порядке, возникает риск гонки состояний, например, на что я и наткнулся -- 
событие завершения задачи обрабатывалось раньше события её создания. Разработать решение этой проблемы мне помог
старший разработчик, предложивший использовать очереди на уровне rootProcessInstanceId. Класс $CamundaEventProcessMap$ хранит словарь, где ключ — это
$rootProcessInstanceId$, а значение — очередь событий для данного процессного экземпляра. При появлении
нового события оно добавляется в очередь своего процесса. Класс $CamundaEventProcessDeque$ представляет очередь событий одного
процессного экземпляра. При создании он сразу запускает асинхронный цикл обработки: поочерёдно извлекает
события из очереди и передаёт их в процессор, дожидаясь завершения каждого через .join().
Консюмер $EventListenerConsumer$ подписывается на топик, при получении каждого сообщения помещает событие в $CamundaEventProcessMap$ и сразу
подтверждает offset. Все запронутые классы и конфигурация прикреплены в \texttt{Приложении Б}.

\subsubsection{Затронутое API}
В процессе разработки плагина я успел поработать с несколькими уровнями API Camunda BPM.
                                                                        
Основной точкой расширения стал интерфейс $ProcessEnginePlugin$ из пакета
org.camunda.bpm.engine.impl.cfg. Он предоставляет три метода          
жизненного цикла: $preInit$, $postInit$ и $postProcessEngineBuild$. Для целей
плагина был использован только $preInit$. Именно через метод
$setHistoryEventHandler$ производится подмена стандартного обработчика.

Интерфейс $HistoryEventHandler$ из пакета
org.camunda.bpm.engine .impl.history.handler определяет два метода:
handleEvent(HistoryEvent) для обработки одиночного события и
handleEvents(List<HistoryEvent>) для пакетной обработки. 
Как раз это достаточно центровое место для интеграции.

Для работы с конкретными типами исторических событий использовались три
класса-сущности из пакета org.camunda.bpm.engine.impl.history.event:
событие элемента активности процесса;
событие пользовательской задачи;
событие экземпляра процесса.
                                                                        
Все они наследуются от базового класса $HistoryEvent$, который    
предоставляет общие поля: тип события,           
$processDefinitionKey$, $processInstanceId$, $rootProcessInstanceId$.        
Конкретные классы добавляют специфические поля — например $taskId$ и
$taskDefinitionKey$, $state$ и $endTime$.

Поля $eventType$ и     
$processDefinitionKey$, оказалось могут быть пустыми при определённых внутренних
событиях движка, поэтому я добавлял валидацию любых входных значений перед любой обработкой.

\subsubsection{Дописывание API}

Помимо использования API Camunda, в рамках проекта я сформировал
API клиентской библиотеки $event-listener-client$, 
предназначенный для сервисов-потребителей.
                                                                        
Интерфейс $EventListenerHandler$ я намеренно сделал максимально
простым: три метода, по одному на каждый класс событий, каждый         
принимает типизированный объект и возвращает объект для асинхронной обработки.

Для подключения клиентской библиотеки к проекту-потребителю достаточно 
добавить зависимость $event-listener-client$ в $pom.xml$ и реализовать
интерфейс $EventListenerHandler$. Библиотека         
автоматически настраивает консюмер, десериализацию и механизм
упорядоченной обработки.

Параметры подключения к Kafka задаются через переменную окружения. 
Если переменная не задана, библиотека
использует значение по умолчанию kafka:9876.

\subsubsection{Результаты сообщений из брокера}

После запуска Camunda с подключённым плагином в топик camunda-events
начинают поступать сообщения — по одному на каждый класс    
событий. Все сообщения сериализуются в JSON.
                                                                        
Сообщение о старте активности (например, сервисной задачи или шлюза)   
имеет вид:

\begin{lstlisting}[
  caption={Сообщение о старте активности},
  label={lst:activity-log}
]
{               
  "eventType": "START",
  "eventClass": "ACTIVITY_INSTANCE",
  "rootProcessInstanceId": "523a0ff1-9c5e-404d-8f98-7a02096a29d7",
  "processInstanceId": "313fab8c-606a-4abd-a3f1-a6c7b027a3d7",                                 
  "processDefinitionKey": "kigimi-process",
  "endTime": null,                                                     
  "activityId": "serviceTask_312d1",
  "activityType": "serviceTask"                                        
}
\end{lstlisting}
                                                                                                                                   
Сообщение о событии пользовательской задачи (создание, завершение или  
удаление):

\begin{lstlisting}[
  caption={Сообщение о событии пользовательской задачи},
  label={lst:activity-log}
]
{               
  "eventType": "CREATE",
  "eventClass": "TASK_INSTANCE",
  "rootProcessInstanceId": "57271022-1de0-4623-8c01-660ae3e743b3",                             
  "processInstanceId": "dbde1229-5675-4745-97fa-61241c2e9b8e",
  "processDefinitionKey": "subfolder-lock",                                
  "endTime": null,
  "taskId": "userTask_32das3",                                            
  "taskDefinitionKey": "userTask_approve"
}  
\end{lstlisting}                                                                    
                
Сообщение о завершении процесса:                                       
\begin{lstlisting}[
  caption={Сообщение о завершении процесса},
  label={lst:activity-log}
]
{                                                                      
  "eventType": "END",
  "eventClass": "PROCESS_INSTANCE",
  "rootProcessInstanceId": "d2fa2c69-4028-46c3-bc7f-191ef75cb895",
  "processInstanceId": "69383101-2a7c-4da6-a44a-cf42fcaecc05",
  "processDefinitionKey": "folder-first-pop-work",                                
  "endTime": "2025-03-09T10:42:17.000+00:00",
  "state": "COMPLETED"                                                 
}   
\end{lstlisting}

Поле $eventClass$ присутствует в каждом сообщении и является ключевым для
десериализации. Поле $rootProcessInstanceId$ позволяет связывать
события дочерних вызванных подпроцессов с корневым процессом — это
поле используется как ключ в $CamundaEventProcessMap$ для
упорядоченной обработки.

\subsubsection{Проверка корректного восстановления}

Для более прозрачной проверки в $EventListenerConsumer$ я добавил лог, полная реализация есть в \texttt{Приложении Б}:

\begin{lstlisting}[
  caption={Логи при падении сервиса},
  label={lst:activity-log}
]
@Service                                                               
  @RequiredArgsConstructor
  public class EventListenerConsumer {

    private static final Logger LOG =                                    
  LoggerFactory.getLogger(EventListenerConsumer.class);
                                                                         
    private final EventListenerHandler eventListenerHandler;
    private final CamundaEventProcessMap map;

    @KafkaListener(topics = KafkaConsumerConfig.EVENT_TOPIC_NAME,        
        containerFactory = "baseEventListenerContainerFactory")
    public void consumeEvent(BaseEvent event, Acknowledgment             
  acknowledgment) {
      LOG.info("Received event: {}", event);
      map.put(event, e -> {
        ...                                                              
      });         
      acknowledgment.acknowledge();                                      
    }             
  }
\end{lstlisting}

Для воспроизведения сценария восстановления после перезапуска я        
попросил развернуть девопсера на тестовом стенде: Camunda с подключённым плагином, Kafka и     
сервис бизнес логкии с потребителем с реализацией $EventListenerHandler$. В Camunda был   
запущен тестовый процесс с несколькими последовательными пользовательскими задачами.
Далее намеренно завершил       
сервис с потребителем в момент обработки событий. Для этого я отправил $kill -9$ процессу. После этого
я перезапустил сервис. При завершении процесса по логике offset не был закоммичен,      
поэтому при следующем старте консюмер начал читать сообщения с
последней подтверждённой позиции. Ниже приведены логи, полученные в    
ходе воспроизведения:
                             
\begin{lstlisting}[
  caption={Логи при падении сервиса},
  label={lst:activity-log}
]
2026-03-24 11:23:41.512  INFO 18432 --- [ntainer#0-0-C-1]
o.s.k.l.ConcurrentMessageListenerContainer : baseEventListenerContainerFactory:
stopping
2026-03-24 11:23:41.513  INFO 18432 --- [ntainer#0-0-C-1]
o.a.k.clients.consumer.KafkaConsumer      : [Consumer
clientId=consumer-rabis-group-1, groupId=rabis-group] Revoke previously assigned
partitions [camunda-events-0]
2026-03-24 11:23:41.514  INFO 18432 --- [ntainer#0-0-C-1]
o.s.k.l.ConcurrentMessageListenerContainer : baseEventListenerContainerFactory:
stopped normally, last committed offset: {camunda-events-0=14}
\end{lstlisting}
                   
                                                                        
В данном случае последний закоммиченный offset равен 14, но я следил за топиком и там было на момент падения уже больше сообщений, что как раз и является не подтвержённым смещением.                                
                
\begin{lstlisting}[
  caption={Логи при перезапуске сервиса},
  label={lst:activity-log}
]
 2026-03-24 11:24:05.301  INFO 19104 --- [           main]
o.s.b.w.embedded.tomcat.TomcatWebServer   : Tomcat initialized with port(s): 8080
  (http)
2026-03-24 11:24:05.744  INFO 19104 --- [ntainer#0-0-C-1]
o.s.k.l.ConcurrentMessageListenerContainer : baseEventListenerContainerFactory:
partitions assigned: [camunda-events-0]
2026-03-24 11:24:05.745  INFO 19104 --- [ntainer#0-0-C-1]
o.a.k.clients.consumer.KafkaConsumer      : [Consumer
clientId=consumer-rabis-group-1, groupId=rabis-group] Setting offset for
partition camunda-events-0 to the committed offset FetchPosition{offset=14,
offsetEpoch=Optional.empty,
currentLeader=LeaderAndEpoch{leader=Optional[kafka:9092], epoch=0}}
2026-03-24 11:24:05.801  INFO 19104 --- [ntainer#0-0-C-1]
r.r.c.e.client.EventListenerConsumer      : Received event:
TaskInstanceEvent(eventType=CREATE, eventClass=TASK_INSTANCE,
rootProcessInstanceId=3c8f1a02-d947-11ef-b17e-0242ac130003,
processInstanceId=3c8f1a02-d947-11ef-b17e-0242ac130003,
processDefinitionKey=contract-approval, endTime=null,
taskId=4d9e2b13-d947-11ef-b17e-0242ac130003, taskDefinitionKey=task_legal_review)
2026-03-24 11:24:05.803  INFO 19104 --- [ntainer#0-0-C-1]
r.r.c.e.client.EventListenerConsumer      : Received event:
TaskInstanceEvent(eventType=COMPLETE, eventClass=TASK_INSTANCE,
rootProcessInstanceId=3c8f1a02-d947-11ef-b17e-0242ac130003,
processInstanceId=3c8f1a02-d947-11ef-b17e-0242ac130003,
processDefinitionKey=contract-approval, endTime=null,
taskId=4d9e2b13-d947-11ef-b17e-0242ac130003, taskDefinitionKey=task_legal_review)
2026-03-24 11:24:05.805  INFO 19104 --- [ntainer#0-0-C-1]
r.r.c.e.client.EventListenerConsumer      : Received event:
ProcessInstanceEvent(eventType=END, eventClass=PROCESS_INSTANCE,
rootProcessInstanceId=3c8f1a02-d947-11ef-b17e-0242ac130003,
processInstanceId=3c8f1a02-d947-11ef-b17e-0242ac130003,
processDefinitionKey=contract-approval, endTime=2026-03-24T11:23:40.991+00:00,
state=COMPLETED)
\end{lstlisting}
                                                                  
Ожидаемое поведение: при жёстком завершении     
сервиса и незакоммиченных offset брокер повторно доставляет          
пропущенные сообщения после перезапуска, а механизм                    
$CamundaEventProcessDeque$ обеспечивает их обработку в корректном
порядке.

\subsubsection{Вывод}
В ходе задачи была внедрена новая подсистема, которая основывается на событийном подходе, который в свою очередь призван убрать опрос персистенции.
Поэтому получении истории и статуса токенов переведено на асинхронную схему, $Camunda$ самостоятельно публикует уведомления о событиях в брокер сообщений. 
На стороне бизнес-логики реализован обработчик подписки на эти уведомления без периодического опроса. Подключение брокера выполнено в виде плагина Camunda Engine, активируемого при старте standalone-сборки. 
Система корректно восстанавливает состояние подписки после перезапуска.
